% res_pair_swing.tex -- DRAFT 23 Oct 2026: the pair's 1.4 Hz swing on the rig (findings; exploratory throughout).
% Sources: rig review docs/reports/rig_review_2026-10-22.md (finding 1, reviewer [weak]: one flight per config, slope
% flights last and h2slope on drained packs, not interleaved); card docs/experiments/2026-10-22_velocity_filter.md (M1:
% the ring normal's horizontal tip band-passed 1.1-1.9 Hz, rms over (a) 0-6 s after lift-off and (b) the hover from the
% top to LAND / drop / detection; tools/pair_swing.py); card docs/experiments/2026-10-23_pair_mode_timing.md (rounds A and
% B, bars frozen 23 Oct 00:42; round C frozen 08:05, R2598-R2601, null; round D reduced, frozen 11:14, R2691-R2694, both
% arms confirm on the thr060 testbed); docs/rig/betaflight_2026-10-22.md (QUAD3
% read-only CLI; ELRS RF mode 8 read as 333 Hz Full from the 14 Oct link statistics: an inference from the mode number).
% Rig rows: RIG-1022-h2, -h2t9p, -h2rts (EMA), -h2slope, -f2slope (slope); -h2m (planned pull) left out (its ring stalled
% 0.37 m short). Twin rows: round A R2449 (cd060 re-fly), R2450/R2451 (thr060), R2452/R2453 (att060), R2454/R2455
% (thr040); round B R2465/R2470 (rt ema 0.55 m), R2471/R2472 (rt slope), R2473/R2474 (rt ema 1 m), R2475/R2476 (rt slope
% 1 m), R2477/R2478 (_rt_norate), R2479/R2480 (_rt050), R2481/R2482 (_rt_nokff). rate_kff 0.027 after the calibration
% run R2456 (0.012) fell outside its band; R2465 is the calibration re-fly, in band at its edge (pitch 35/37 ms).
% Two runner FAILs in round A (R2449, R2453): an envelope fault at touchdown (77.9 deg), after the scored window.
% Decisions 23 Oct late morning: (2) system_overview's latency corrected (done, outside these drafts); (3) round C null,
% round D reduced to the config-only rig candidates on the thr060 testbed; Wed 28's 1 m rig pair is the test; (4)
% rate_kff, thr_tau_s, the body-rate filter emulation and the tip-hinge world DELETED (11:19; backups in the scratchpad's
% del_rt/); the per-channel command delay kept, default off.
% NUMBERS DIFFER (for Wesley): system_overview.tex (i) now gives body rate 17-24 (rig) against 49-57 ms (twin) and
% throttle-to-lift ~15-25 ms in the twin, from the Betaflight note's design numbers; Table pair_timing below gives the
% card's later loop_timing.py readings, 25-35 against 48-62 ms and 17-29 ms. Not edited outside these drafts.

\subsection{The pair's swing on the rig}
\label{sec:res_pair_swing}

This section records a rig finding that the twin does not yet explain. Two UAVs carrying the ring on opposite plates,
the stage that every chain of releases passes through, swing at about 1.4\,Hz on the rig, and the twin's pair does
not. It asks which command path can carry such a swing, and whether a twin given the rig's measured timing reproduces
it. Everything here is exploratory; the first pair on 1\,m cables, on the rig day of 28 October, is the test.

\subsubsection{On the rig}

On 22 October, with the motion-capture velocity filter that was then the rig's default (an exponential moving average),
every two-UAV lift on 0.55\,m rods swung at 1.3--1.6\,Hz, and the swing grew (Table~\ref{tab:pair_rig},
Figure~\ref{fig:pair_rig}). One flight lost the ring when a magnet let go near the top of the lift; the others ended
with the operator's LAND. A higher throttle cap, other plates and references shifted by the plan's age left the swing
as it was. With a filter about 30\,ms faster (the slope of a short window) on the two UAVs the swing halved and
decayed, and the pair carried the ring through a whole figure-eight, the first full lap with two UAVs, with the ring
hanging tilted about the pair's line. Each configuration flew once, and the two slope flights came last, one on drained
packs: the filter is a lead, not a result, until the interleaved repeat of the next rig day.

\begin{table}[htbp]
\centering
\footnotesize
\setlength{\tabcolsep}{3pt}
\caption[The pair on the rig]{Two UAVs carrying the 0.86\,kg ring on opposite plates on the rig, 0.55\,m rods, hover
unless stated, 22 October; one flight each. Velocity filter: the motion-capture velocity estimate the trackers use,
an exponential moving average (about 45\,ms of lag) or the slope of a short window (about 17\,ms). Swing: the rms of
the ring's tip band-passed at 1.1--1.9\,Hz, over the first 6\,s after lift-off (climb) and over the hover. A flight on
the planner's pull, whose ring stalled short of its height, is left out. $^a$\,The flight hovered less than a second.}
\label{tab:pair_rig}
\begin{tabularx}{\textwidth}{c>{\raggedright\arraybackslash}p{1.6cm}>{\raggedright\arraybackslash}p{3.4cm}cc>{\raggedright\arraybackslash}X}
\toprule
 & Velocity & & \multicolumn{2}{c}{Swing ($^\circ$)} & \\
\cmidrule(lr){4-5}
Flight & filter & Change & Climb & Hover & Outcome \\
\midrule
1 & average & plates 0 and 6 & 3.72 & --$^a$ & a magnet let go near the top of the lift (6.4\,N peak pull); ring dropped, every UAV landed \\
2 & average & plates 11 and 5; throttle cap 0.9 & 2.60 & --$^a$ & operator LAND, swing growing \\
3 & average & as 2; references shifted by the plan's age & 2.99 & --$^a$ & operator LAND, swing growing \\
4 & slope & as 2 & 1.46 & 0.79 & swing decaying; landed \\
5 & slope & as 2; figure-eight at 0.2\,m/s & 1.35 & 0.74 & the whole lap; the ring hung about 7$^\circ$ about the pair's line \\
\bottomrule
\end{tabularx}
\end{table}
% 1 RIG-1022-h2 (1.41 Hz; QUAD3's magnet at ring z 0.94 m, lift 95 %; UNANNOUNCED DETACH, DROP AND LAND, ring 62 deg);
% 2 -h2t9p (1.49 Hz; throttle never above 0.77, so the cap was never used; LAND 6.2 s after lift-off); 3 -h2rts (1.45 Hz,
% ref_time_shift on); 4 -h2slope (the 0.71 Hz pendulum left); 5 -f2slope (3-D RMSE 3.7 cm, tilt peak 14.3, mean on the
% path 7.3 deg, lag 0.10 s). Climb values: the velocity-filter reviewer's, reproduced by tools/pair_swing.py within 0.1 deg.

\pendingfig{fig:pair_rig}{The pair's swing on the rig, 22 October: the ring's tip band-passed at 1.1--1.9\,Hz against
time from lift-off, two UAVs on opposite plates, 0.55\,m rods. (a) The three flights with the averaging velocity
filter, the swing growing until LAND or a magnet's release; (b) the two flights with the slope filter, the swing
decaying.}{the five rig flights of Table~\ref{tab:pair_rig}}
% Data: results/rig/2026-10-22/{h2,h2t9p,h2rts,h2slope,f2slope}_logs (ring pose). Tool: tools/pair_swing.py (its
% band-pass and windows) -> a new thesis_rig_figures function through the thesis-figure skill.

\subsubsection{Hardware and timing}

The flight controllers run Betaflight in rate mode: each holds the three body rates it is sent, with a feed-forward on
the rate command and a fast rate loop, and passes the collective throttle to the motors with no loop on it. The tracker
sends its commands at 50\,Hz, the laptop writes the latest frame to each radio every 10\,ms, and the ExpressLRS link
runs at 333\,Hz (inferred from the mode its link statistics report). Radio and flight controller together add only about 10--15\,ms
(Table~\ref{tab:pair_timing}). Measured end to end, the rig differs from the twin per channel: it turns a body-rate
command into body rate faster than the twin, and a throttle command into lift slower, since nothing on the flight
controller closes a loop on thrust and the motors' spin-up shows in full; and the body rate the trackers see lags by
the motion-capture node's own filter. A single command delay in the twin therefore stood for two different things.
No setting of the flight controllers was changed for this work.

\begin{table}[htbp]
\centering
\footnotesize
\caption[The rig's command timing]{The rig's command timing, from the flight controller's configuration (read only)
and the rig's logs of 22 October, against the twin before the timing rounds. Estimates per stage add to the radio and
flight controller's share; the end-to-end paths are measured as an equivalent delay. $^a$\,From the trackers' log
(14--28\,ms from the recordings). $^b$\,The motion-capture node's moving average on the body rates (always on).}
\label{tab:pair_timing}
\begin{tabular}{lcc}
\toprule
 & Rig (ms) & Twin (ms) \\
\midrule
\multicolumn{3}{l}{\emph{Stages, estimated}} \\
Laptop frame hold (a 100\,Hz writer after the 50\,Hz tracker) & 0--10 & -- \\
Radio packet wait and air time at 333\,Hz & 2--3 & -- \\
Command smoothing on the flight controller & 2 (6 at worst) & -- \\
Gyro and rate-loop filters & 1--2 & -- \\
Radio and flight controller together & 10--15 & -- \\
\midrule
\multicolumn{3}{l}{\emph{Paths, measured end to end}} \\
Body-rate command to body rate & 25--35$^a$ & 48--62 \\
Throttle to lift & 44--67 & 17--29 \\
Body rate as the trackers see it$^b$ & 36--48 & 4--5 \\
\bottomrule
\end{tabular}
\end{table}
% Betaflight 4.3.1 (STM32F722), rate profile ACTUAL 100 deg/s linear; rate PID roll P45 I80 D40 F120, pitch P47 I84 D46
% F125, yaw F120; gyro and PID 8 kHz; RC smoothing auto, factor 30; DShot600, airmode; no RPM filter (bidirectional DShot
% off); failsafe: drop after 1.5 s. Twin paths: R2381/R2384/R2385/R2388 after their 60 ms delay was removed (tools/
% loop_timing.py att / lift / filter). QUAD3 read; QUAD1, 2 and 4 assumed the same (not read).

\subsubsection{In the twin}

Four twin rounds asked which path carries the swing, whether the rig's timing or its hinge reproduces it, and what
removes it, each on the twin's pair hover with its delays switched on after the lift (Table~\ref{tab:pair_twin}). In
the first, a 60\,ms delay on every channel made the twin's pair swing; the same delay on the throttle alone swung as
much, and on the body rates alone the pair stayed calm. In the twin, then, the throttle path carries the swing. With
the delay on the body rates alone, though, a slower mode grew until LAND whose shape is the rig's: the ring slid along
the pair's line, and the thrust along the line fed the swing (Figure~\ref{fig:pair_mode}).

In the second round the twin was given the rig's measured timing per channel: the throttle 35\,ms slower, a
feed-forward on the rate loop set so that the body-rate path matches the rig's, and the motion-capture node's filter
on the body rates. It stayed calm with either velocity filter, on 0.55\,m rods and on 1\,m cables, and without the
filter or the feed-forward, which falsifies the reading that the rig's timing alone makes the rig's pair swing. With
the throttle 50\,ms slower it swung, but in the twin's own mode, the UAVs' heights moving with the plates, not in the
rig's. For the first 1\,m pair on the rig the twin therefore predicts no swing with either filter, a weak prediction
from a twin that missed the 0.55\,m pair.

The third round moved the twin's hinge to the magnet tip, 5.75\,cm above the ring's centre of mass as measured on the
rig, with the rig's timing and the planner typed either as on the rig or to the hinge. The pair stayed calm both ways,
so in the twin the hinge's height does not make the rig's swing either. With the second round's reading falsified and
the third null, the rig-timing emulation and the tip-hinge world were removed from the twin; the per-channel delay
stays, off by default, as the instrument of the first and fourth rounds.

The fourth round tried, on the swinging twin, the two changes the rig can fly without new code: the slope filter, and
the trackers on the planner's pull instead of the measured one. Its testbed was the first round's 60\,ms on the
throttle alone. Both removed the swing; on the planner's pull the swing's mode was gone altogether, which places the
twin's swing in the measured pull. The testbed's swing is the twin's own, the UAVs' heights moving with the plates,
not the rig's, so the round shows that each change damps a throttle-delayed swing in the twin, not that either removes
the rig's. The slope filter is already the pair's comparison on the rig on 28 October; the planner's pull would be a
change for the pair alone, since the measured pull is what keeps a release's tilt small (Section~\ref{sec:res_leave}).

\begin{table}[htbp]
\centering
\footnotesize
\setlength{\tabcolsep}{3pt}
\caption[The pair's swing in the twin]{The pair's swing in the rig twin: two UAVs on opposite plates, the 0.86\,kg
ring, hover, 0.55\,m rods unless stated, the averaging velocity filter and the measured pull unless stated; each delay
switched on after the lift. Swing: the hover swing of Table~\ref{tab:pair_rig}; a twin pair counts as swinging at
1.0$^\circ$ or more. Rig timing: the throttle 35\,ms slower, a rate feed-forward matching the rig's body-rate path, and
the motion-capture node's filter on the body rates. Round 3: the rig's timing with the hinge at the magnet tip. Round
4: the first round's 60\,ms on the throttle only, its swing 7.36--7.39$^\circ$ without a change.}
\label{tab:pair_twin}
\begin{tabularx}{\textwidth}{c>{\raggedright\arraybackslash}Xccl}
\toprule
Round & Arm & $n$ & Swing ($^\circ$) & Reading \\
\midrule
1 & 60\,ms on every channel & 1 & 6.95 & swings \\
 & 60\,ms on the throttle only & 2 & 7.36--7.39 & swings as with every channel \\
 & 60\,ms on the body rates only & 2 & 0.56--0.81 & calm; a slower mode grows \\
 & 40\,ms on the throttle only & 2 & 0.10--0.12 & calm \\
\midrule
2 & rig timing & 2 & 0.066--0.081 & calm \\
 & rig timing, slope filter & 2 & 0.029 & calm \\
 & rig timing, 1\,m cables & 2 & 0.065--0.069 & calm \\
 & rig timing, slope filter, 1\,m cables & 2 & 0.029 & calm \\
 & rig timing without the body-rate filter & 2 & 0.077--0.082 & calm \\
 & rig timing without the rate feed-forward & 2 & 0.070--0.089 & calm \\
 & rig timing, the throttle 50\,ms slower & 2 & 4.07--5.14 & swings, in the twin's own mode \\
\midrule
3 & planner typed as on the rig & 2 & 0.08--0.09 & calm \\
 & planner typed to the hinge & 2 & 0.08 & calm \\
\midrule
4 & slope filter & 2 & 0.06--0.07 & calm \\
 & the planner's pull & 2 & 0.01 & calm; no mode at the swing's frequency \\
\bottomrule
\end{tabularx}
\end{table}
% Round A confirms (frozen bars: thr060 >= 0.5x the re-fly, att060 < 0.3x); att060's slower mode 1.0-1.1 Hz, ring slide
% 0.64-0.67 cm/deg, c_u < 0. Round B falsifies (both ema runs < 1 deg and both ring slides < 0.1 cm/deg; items 1, 4, 5
% fail). _rt050 grows to ~19 deg p-p, c_u > 0 (the twin's vertical mode). Checks: att pitch 31-39 ms (rig 25-35).
% Round C (R2598/R2600 typed0175, R2599/R2601 typed0575): the testbed bar fails (typed0175 M1(b) 0.09/0.08 < 1.0 deg;
% ring slide 0.054/0.050 < 0.15): null, no margin claim. Round D (R2691/R2692 slope, R2693/R2694 cable_source model; no-fix
% arm R2450/R2451 7.39/7.36): both arms CONFIRM (both runs <= 3.69 deg, no new fault); D2's f0 1.00 Hz, no 1.5 Hz mode;
% c_u > 0 (the twin's mode). D2 is a candidate for the pair only, after Wed's slope A/B, not a default (the measured pull
% carries Tier A's gains, R1706-R1711). D3-D5 not flown. rate_kff / thr_tau_s / the rate filter / the tip world deleted
% 23 Oct 11:19 (card section 8).

\pendingfig{fig:pair_mode}{The shape of the pair's swing at its frequency, per degree of the ring's see-saw: how far
the ring slides along the pair's line against how far the UAVs' heights move. The rig's flights with the averaging
filter and with the slope filter, against the twin's delayed arms (60\,ms on every channel, on the body rates only)
and the rig-timing twin; the rig's swing slides the ring, the twin's moves the UAVs' heights.}{the rig flights of
Table~\ref{tab:pair_rig} and the twin arms of Table~\ref{tab:pair_twin}, rounds 1--4}
% Data: tools/pair_mode.py on RIG-1022-h2, -h2t9p, -h2rts, -h2slope, -f2slope and on R2381/R2384/R2385/R2388 (cd060),
% R2452/R2453 (att060), R2465/R2470 (rt), R2598-R2601 (round C), R2691-R2694 (round D); slide and dz in cm per deg, c_u as
% the marker's fill (feeds or
% damps). Tool: tools/pair_mode.py -> a new figure function through the thesis-figure skill.

\paragraph{Answer.} The pair's 1.4\,Hz swing is a rig finding. In the twin a delayed throttle carries a swing of the
same frequency, but a twin with the rig's measured timing stays calm, and so does one with the rig's hinge height, so
neither explains the rig's swing. On the twin's own throttle-delayed swing, the faster velocity filter and the
planner's pull each remove it, which says what damps the twin's mode, not the rig's. The faster filter, which calmed
the rig's pair in one flight, waits for its interleaved repeat on the rig, with the first pair on 1\,m cables.
